% Part: propositional-logic
% Chapter: propositional-logic 
% Section: completeness

\documentclass[../../../include/open-logic-section]{subfiles}

\begin{document}

\olfileid{pl}{prp}{com}

\olsection{کامل‌بودن منطق گزاره‌ای}

\begin{defn}
\ollabel{def:MaxCon}
مجموعهٔ~$\Gamma$ از~!!{formula}s را \emph{بیشینه‌سازگار} می‌نامیم اگر
سازگار باشد و هرگاه~$\Delta$ مجموعه‌ای سازگار و
$\Gamma\subseteq\Delta$ باشد، آنگاه~$\Gamma=\Delta$.
\end{defn}

\begin{lem}[لم صدق]
  \ollabel{lem:truth}
فرض کنید~$\Gamma$ بیشینه‌سازگار باشد؛ آنگاه:
\begin{enumerate}
\item \ollabel{lem:truth-1} $!A\in\Gamma$ اگر و تنها اگر
  $\lnot !A\notin\Gamma$؛
\item \ollabel{lem:truth-2} $\Gamma\Proves !A$ اگر و تنها اگر
  $!A\in\Gamma$؛
\item \ollabel{lem:truth-3} $!A\lif !B\in\Gamma$ اگر و تنها اگر
  یا~$!A\notin\Gamma$ یا~$!B\in\Gamma$.
\end{enumerate}
\end{lem}

\begin{proof}
\begin{enumerate}
\item فرض کنید~$!A\in\Gamma$. سازگاری نتیجه می‌دهد که
  $\lnot !A\notin\Gamma$. برعکس، فرض کنید~$\lnot !A\notin\Gamma$ و
  $!A\notin\Gamma$. در این صورت هیچ‌یک در~$\Gamma$ نیست؛ بنا بر
  \olref[axd]{prop:phi} یکی از~$\Gamma\cup\{!A\}$ و
  $\Gamma\cup\{\lnot !A\}$ امتدادی درست و سازگار از~$\Gamma$ است، که
  با بیشینگی آن تناقض دارد. پس~$!A\in\Gamma$.
\item اگر~$!A\in\Gamma$، آشکار است که~$\Gamma\Proves !A$. از سوی
  دیگر، فرض کنید~$\Gamma\Proves !A$. اگر~$!A\notin\Gamma$ باشد، بنا بر
  \olref{lem:truth-1} داریم~$\lnot !A\in\Gamma$. پس هم
  $\Gamma\Proves !A$ و هم~$\Gamma\Proves\lnot !A$، که تناقض است.
\item تمرین.
\end{enumerate}
\end{proof}

جهت چپ‌به‌راستِ~\olref{lem:truth}\olref{lem:truth-2} نشان می‌دهد که
$\Gamma$ \emph{استنتاجاً بسته} است؛ یعنی اگر~$\Gamma\Proves !A$، آنگاه
$!A\in\Gamma$، برای هر~$!A$.

\begin{prob}
  برهان~\olref{lem:truth} را کامل کنید.
\end{prob}

\begin{thm}[کامل‌بودن]
\ollabel{thm:completeness}
اگر~$\Gamma$ سازگار باشد، ارضاپذیر است.
\end{thm}

\begin{proof}
فرض کنید~$!A_0,!A_1,\dots$ فهرستی فراگیر از همهٔ~!!{formula}s زبان
باشد. دنباله‌ای افزایشی از مجموعه‌های~!!{formula}s، یعنی
$\Gamma_0,\Gamma_1,\dots$، را با قرار دادن موارد زیر به‌طور بازگشتی
تعریف می‌کنیم:
\begin{align*}
\Gamma_0 & = \Gamma\\
\Gamma_{n+1} &  =
\begin{cases}
  \Gamma_n \cup \{!A_n\} & \text{اگر $\Gamma_n\cup\{!A_n\}$
    سازگار باشد؛}\\
  \Gamma_n \cup \{\lnot !A_n\} & \text{در غیر این صورت.}
\end{cases}
\end{align*}
سپس تعریف می‌کنیم:
\[
\Gamma^*=\bigcup_{0\le n}\Gamma_n.
\]
اکنون برهان با اثبات پیاپیِ واقعیت‌های زیر ادامه می‌یابد:
\begin{enumerate}
\item برای هر~$n$، مجموعهٔ~$\Gamma_n$ سازگار است (با استقرا بر~$n$ و
  استفاده از~\olref[axd]{prop:phi})؛
\item $\Gamma^*$ سازگار است؛
\item $\Gamma^*$ بیشینه است.
\end{enumerate}
سپس ارزش‌گذاری~$v$ را با~$v(p_i)=\True$ اگر و تنها اگر
$p_i\in\Gamma^*$ تعریف می‌کنیم. با استقرا بر~$!A$ نشان می‌دهیم که برای
!!{sentence}s پیچیده‌تر نیز عضویت در~$\Gamma^*$ با صدق بر حسب~$v$
منطبق است: $\overline{v}(!A)=\True$ اگر و تنها اگر
$!A\in\Gamma^*$. به‌ویژه، $v$ مجموعهٔ~$\Gamma^*$ را ارضا می‌کند؛ و چون
$\Gamma\subseteq\Gamma^*$، $\Gamma$ نیز ارضاپذیر است، چنان‌که می‌خواستیم.
\end{proof}

\begin{cor}
اگر~$\Gamma\Entails !A$، آنگاه~$\Gamma\Proves !A$.
\end{cor}

\begin{proof}
اگر~$\Gamma\Proves/ !A$، بنا بر~\olref[axd]{prop:prov-incons} مجموعهٔ
$\Gamma\cup\{\lnot !A\}$ سازگار است. پس به موجب قضیهٔ کامل‌بودن،
$\Gamma\cup\{\lnot !A\}$ ارضاپذیر است و~$\Gamma\Entails/ !A$.
\end{proof}

\begin{prop}[قضیهٔ فشردگی]
\ollabel{prop:compactness}
$\Gamma$ ارضاپذیر است اگر و تنها اگر هر زیرمجموعهٔ \emph{متناهی}
$\Gamma_0$ از~$\Gamma$ ارضاپذیر باشد.
\end{prop}

\begin{proof}
$\Gamma$ ارضاناپذیر است اگر و تنها اگر ناسازگار باشد؛ اگر و تنها اگر
زیرمجموعه‌ای متناهی چون~$\Gamma_0$ از~$\Gamma$ ناسازگار باشد؛ و اگر و
تنها اگر زیرمجموعه‌ای متناهی چون~$\Gamma_0$ از~$\Gamma$ ارضاناپذیر باشد.
\end{proof}

\begin{cor}
$\Gamma\Entails !A$ اگر و تنها اگر زیرمجموعه‌ای متناهی چون~$\Gamma_0$
از~$\Gamma$ وجود داشته باشد که~$\Gamma_0\Entails !A$.
\end{cor}

\end{document}
